% Generated draft: requires mathematical and source review.
\documentclass[11pt]{article}
\usepackage[margin=25mm]{geometry}
\usepackage{fontspec}
\setmainfont{Latin Modern Roman}
\usepackage{amsmath,amssymb,mathtools,unicode-math}
\setmathfont{Latin Modern Math}
\usepackage{enumitem,needspace,xurl,hyperref}
\setlength{\parindent}{0pt}
\setlength{\parskip}{5pt}
\setlength{\emergencystretch}{3em}
\newcommand{\N}{\mathbb{N}}
\newcommand{\Z}{\mathbb{Z}}
\newcommand{\Q}{\mathbb{Q}}
\newcommand{\R}{\mathbb{R}}
\newcommand{\C}{\mathbb{C}}
\newcommand{\F}{\mathbb{F}}
\newcommand{\A}{\mathbb{A}}
\newcommand{\PP}{\mathbb{P}}
\newcommand{\E}{\mathbb{E}}
\newcommand{\eps}{\varepsilon}
\newcommand{\dd}{\,\mathrm{d}}
\begin{document}
\section*{UnsolvedMath ID: 801082}
\textbf{Source problem code:} catalogue entry
\subsection{Definitions and mathematical statement}
Let \(f\colon X\to X\) denote the H\'enon dynamical system and phase space intended in the source. The supplied excerpt does not specify the class of maps or the phase space \(X\). For a potential \(\phi\), write its Birkhoff sum as
\[
S_n\phi(x)=\sum_{j=0}^{n-1}\phi(f^j(x)).
\]
Let \(\delta_x\) be the unit point mass at \(x\), and let \(P(\phi)\) denote the pressure intended by the source. An equilibrium state is an invariant Borel probability measure attaining the pressure through the entropy-plus-integral expression
\[
h_{\mu_\phi}(f)+\int_X\phi\,\dd\mu_\phi=P(\phi),
\]
with the admissible measures, integrability conditions, and pressure convention understood relative to the intended system. Those choices are not supplied.

Write \(SP_n\) for the source's collection of saddle periodic points. Under the usual differentiable convention, a saddle periodic point has a hyperbolic derivative of its return map with both contracting and expanding directions. The excerpt does not define \(SP_n\), in particular whether it includes all saddle points fixed by \(f^n\) or only those of least period \(n\).

The source asks: suppose \(\phi\) is sufficiently regular that its equilibrium state \(\mu_\phi\) exists and is unique. Is it true that, as \(n\to\infty\),
\[
\nu_n:=e^{-nP(\phi)}\sum_{x\in SP_n}e^{S_n\phi(x)}\delta_x
\longrightarrow\mu_\phi?
\]
The natural measure-theoretic interpretation is weak convergence of finite Borel measures, namely
\[
e^{-nP(\phi)}\sum_{x\in SP_n}e^{S_n\phi(x)}g(x)
\longrightarrow\int_X g\,\dd\mu_\phi
\]
for every bounded continuous real-valued function \(g\) on \(X\). This topology is an interpretation requiring confirmation, since the excerpt does not specify it. The regularity class of \(\phi\), its scalar range, the hypotheses on \(f\), and the conditions ensuring that these weighted measures are finite are also unspecified. Thus this is a source-relative formulation, not a fully specified assertion for every H\'enon map.
\subsection{Short English statement}
For the H\'enon system intended by the source, does the weighted sum of point masses at saddle periodic points, normalized by \(e^{nP(\phi)}\), converge to the unique equilibrium state of \(\phi\)? The map class, phase space, potential regularity, pressure convention, periodic-point convention, and convergence topology require clarification.
\subsection{Sources}
\begin{itemize}
\item Local UnsolvedMath registry, permanent ID \(801082\), source problem code \texttt{catalogue entry}, title ``Henon Maps: Open Problems --- Question 34.'' The supplied record attributes the problem to Berger and Bedford, with proposed year 2023, and labels it open. This label is supplied registry metadata, not an independently verified mathematical or literature assessment.
\item Berger and Bedford, \emph{Henon Maps: A List of Open Problems} (2023), Question 34. Supplied source URL: \url{https://arxiv.org/abs/2312.03907}. The supplied provenance records access on 2026-10-16 and extraction by \texttt{arxiv\_latex}; that access report has not been independently verified here.
\end{itemize}
\end{document}
